% =====================================================================
%  PAPER OUTLINE — sc-crop + contrast-agnostic v4
%  Working title:
%    "Detect-then-segment: a contrast-agnostic spinal cord cropping
%     pipeline that improves segmentation accuracy and inference speed"
%
%  HOW TO USE THIS FILE
%  --------------------
%  This is a STRUCTURED SKELETON, not a finished paper. Every section is a
%  bullet-point plan. Markers tell you exactly what to write/insert:
%     [WRITE]        -> prose you still need to write
%     [INSERT-RESULT]-> a number/figure/table that ALREADY EXISTS (source given)
%     [FUTURE]       -> a placeholder reserved for upcoming models/experiments
%     [DECISION]     -> an authoring choice to settle before submission
%  It compiles as-is (pdflatex). Replace bullets with prose section by section.
%
%  TARGET VENUE [DECISION]: methods/tool paper. The lab publishes in
%  Medical Image Analysis and Imaging Neuroscience -> assume that audience
%  (medical imaging + neuro), single-column or journal template later.
% =====================================================================

\documentclass[11pt]{article}
\usepackage[utf8]{inputenc}
\usepackage{graphicx}
\usepackage{booktabs}
\usepackage{amsmath}
\usepackage{xcolor}
\usepackage[colorlinks=true,linkcolor=blue,citecolor=blue]{hyperref}
\usepackage{geometry}
\geometry{margin=1in}

% --- authoring markers (render in color so they stand out in drafts) ---
\newcommand{\WRITE}[1]{\textcolor{blue}{\textbf{[WRITE]} #1}}
\newcommand{\RESULT}[1]{\textcolor{teal}{\textbf{[INSERT-RESULT]} #1}}
\newcommand{\FUTURE}[1]{\textcolor{orange}{\textbf{[FUTURE]} #1}}
\newcommand{\DECISION}[1]{\textcolor{red}{\textbf{[DECISION]} #1}}

\title{Detect-then-segment: a contrast-agnostic spinal cord cropping
pipeline that improves segmentation accuracy and inference speed
\\ \large \DECISION{title is a placeholder — settle once results are final}}
\author{Quentin Revillon \and \FUTURE{co-authors}}
\date{Draft outline}

\begin{document}
\maketitle

% =====================================================================
\section*{Reviewer's map of contributions (delete before submission)}
% =====================================================================
% Keep this crystal clear in your head while writing — the paper has THREE
% deliverables. Do not let them blur together:
\begin{itemize}
  \item \textbf{C1 — sc-crop}: a generic, contrast- and modality-agnostic
        spinal-cord \emph{detection + cropping} tool (2D axial YOLO detector
        + connected-component classifier regularization), released as a CLI
        and a Python package, ONNX (CPU) and PyTorch (GPU) backends.
        \emph{This is the methodological core.}
  \item \textbf{C2 — crop ablation}: the same contrast-agnostic nnUNet trained
        on cropped vs full volumes (identical data, labels, split), used to
        \emph{quantify} the benefit of detect-then-segment (Dice + speed) with the
        crop as the only difference.
  \item \textbf{C3 — deployment + monitoring}: integration into Spinal Cord
        Toolbox (\texttt{sct\_deepseg}) and an automated lifelong-learning
        pipeline that tracks morphometric (CSA) drift across model releases.
\end{itemize}
\DECISION{Lead with C1 (the tool generalizes beyond this one model) and use
C2 as the validation. Alternative framing — lead with C2 (a better SC model)
— is weaker because the novelty is the pipeline, not the nnUNet retrain.}

% =====================================================================
\section*{Abstract}\label{sec:abstract}
% =====================================================================
% ~200--250 words. Prose draft below. Numbers in \WRITE{} are NOT yet measured
% and each is gated on a specific experiment — see "EXPERIMENTS REQUIRED" block.
%
% FRAMING (do not drift): the argument is a clean chain --- (1) nnUNet is the
% SOTA for SC segmentation but slow on CPU (clinical bottleneck); (2) two ways to
% speed it up: change the architecture (accuracy cost) OR shrink the FOV by
% cropping (architecture untouched -> accuracy kept); (3) we take route 2 and
% build a fast, robust, REUSABLE, contrast-agnostic CPU crop. The contribution is
% that reusable crop, validated by retraining nnUNet on crops. Accuracy is
% reported as MAINTAINED / non-inferior; the small Dice gain is attributed to
% cropping ONLY via the controlled no-crop ablation (Dataset7001) --- the
% v3-vs-v4 number alone is a confounded deployment comparison.

Spinal cord (SC) segmentation has reached high accuracy across contrasts and
pathologies with the nnUNet-based contrast-agnostic model deployed in the Spinal
Cord Toolbox (SCT). The remaining barrier to clinical adoption is speed. On CPU,
whole-volume inference takes up to $\sim$2\,min \WRITE{[confirm, EXP-SPEED]}.
nnUNet processes the entire volume, yet the cord occupies only a small fraction
of it --- so cropping to a cord-centred region of interest would reduce inference
cost by the same factor, with the architecture left untouched. No standalone,
contrast-agnostic cropping tool exists: existing localisers are either absent
from the pipeline \cite{bedard2025media} or tightly coupled to a specific
segmentation model \cite{lemay2021tumor, gros2021tumor}.

We present \textbf{sc-crop}, a fast, reusable, modality-agnostic detect-and-crop
primitive for spinal cord imaging. A 2D axial YOLO detector runs per slice; a
connected-component classifier removes superior false positives (brain stem, neck
structures), yielding a tight 3D bounding box in milliseconds on CPU. sc-crop is
decoupled from any segmentation model and exposes atomic
\texttt{detect}/\texttt{crop}/\texttt{uncrop} operations, so it embeds unchanged
into nnUNet preprocessing, SCT, or any custom pipeline. Cropping reduces the
field of view by $\eta\approx\WRITE{[5\text{--}15,~\text{EXP-FOV}]}$, cutting
per-volume inference time by \WRITE{[$\times$X, EXP-SPEED]} on CPU. Accuracy is
not sacrificed --- it improves. In a controlled ablation that retrains the same
contrast-agnostic nnUNet on cropped versus full volumes (identical data, labels
and train/val/test split, so that only the crop differs), cropping yields a small
but highly significant Dice gain on a frozen, multi-site, multi-contrast held-out
set ($0.958\pm0.020$ vs.\ $0.955\pm0.028$, $n=510$, 9 contrasts, 14 datasets,
paired Wilcoxon $p=1\times10^{-14}$) and, more importantly, improves robustness:
it lowers variance ($\sigma=0.020$ vs.\ $0.028$) and removes catastrophic failures
(worst-case Dice $0.77$ vs.\ $0.51$). No ground-truth voxel is lost in any subject
($510/510$). sc-crop is open-source,
ONNX-deployable with no deep-learning runtime dependency, and integrated into
SCT.

% ---------------------------------------------------------------------
% EXPERIMENTS REQUIRED to back each abstract claim (NONE optional for those
% claims — without them the sentence must be softened or cut):
%
%   CLAIM "cuts inference time by Xx on CPU"  [\WRITE in abstract]
%     -> EXP-SPEED (MISSING). Like-for-like wall-clock, SAME PyTorch backend:
%        sct_deepseg (full volume) vs sc-segment-pt (crop+nnUNet), same machine,
%        CPU and GPU, n volumes spanning FOV sizes; split detect/segment/total.
%        Cleanest: SAME nnUNet weights on full vs cropped input (ties to EXP-ABL).
%        NOTE: no timing exists anywhere yet (not in metrics.json, not in
%        bench_preprocessing.py which times only YOLO preprocessing). Until run,
%        the "Xx" placeholder is UNSUPPORTED.
%
%   CLAIM "shrinks FOV by eta ~ 5-15"          [\WRITE in abstract]
%     -> EXP-FOV (MISSING but cheap). Compute cropped-bbox-volume / original-
%        volume per test subject from the saved bboxes (NOT voxels_before/after
%        in metrics.json -- those are LABEL voxels, always equal => 0 lost).
%        Report mean + range. This is the confound-free efficiency number.
%
%   CLAIM "at no cost in accuracy / matches-or-exceeds baseline"  [numbers real]
%     -> Already supported by metrics.json (Dice 0.9578 vs 0.9535, n=510,
%        Wilcoxon p=2.2e-25, 70% win). Keep as NON-INFERIORITY, not "crop helps".
%
%   CLAIM "preserves every GT voxel (510/510)"  [real]
%     -> metrics.json aggregate n_bad_crops=0. Solid.
%
%   ENABLES any causal "cropping is responsible" statement (Discussion, NOT
%   abstract) -> EXP-ABL (STAGED, NOT RUN): Dataset7001 no-crop twin
%   (train_contrast_agnostic_nocrop.sh; fill REPLACE_WITH_THE_V4_DATALISTS_FOLDER
%   with /home/quentinr/datalists/<the v4 folder>). Same split/labels/nnUNet,
%   only the crop differs. This is the single biggest reviewer risk: without it
%   the v3-vs-v4 Dice gain is a confounded DEPLOYMENT comparison and cannot be
%   attributed to cropping. The abstract sidesteps this by claiming
%   non-inferiority, not improvement -- keep it that way unless EXP-ABL is run.
% ---------------------------------------------------------------------

% \FUTURE{When new models land, the Dice/speed sentence gets a third comparison
% point — keep the "matches or slightly exceeds" structure so it extends.}
% \DECISION{Scope claim: sc-crop detects the CORD. It accelerates SC + intra-
% medullary tasks (lesions/tumours, GM) whose target lies INSIDE the cord ROI.
% Do NOT claim it accelerates general vertebral-column / whole-spine segmentation
% — a cord-centred box + small margin does not contain the vertebral bodies or
% posterior elements; a reviewer will reject that. Mention whole-spine only as
% outlook, not as a demonstrated benefit.}

% =====================================================================
\section{Introduction}\label{sec:intro}
% =====================================================================
\paragraph{State of the art: accurate but slow.}
Deep learning now dominates spinal cord (SC) segmentation. Across tasks ---
segmenting the cord itself, multiple-sclerosis (MS) lesions, or intramedullary
tumours --- the self-configuring nnUNet framework \cite{isensee2021nnunet} has
become the de facto standard \cite{}, and in the Spinal Cord Toolbox (SCT) the
contrast-agnostic model \texttt{sct\_deepseg spinalcord} \cite{bedard2025media} is
built on it. The field has converged on a single \emph{contrast-agnostic} model
because clinical scans span many contrasts (T1w, T2w, T2$^\ast$, MT, DWI, PSIR,
STIR, MP2RAGE) and per-contrast models do not scale \cite{bedard2025media}. These
models are highly accurate --- but slow: on the CPUs available in most clinical
and many research settings, the deployed model can take up to $\sim$2\,min per
volume \WRITE{[confirm with EXP-SPEED]}, too slow for use at the point of care.
The cause is structural. In a spinal MRI the cord fills only a small fraction of
the field of view (FOV) --- a few tenths of a percent of voxels, and far less for
a lesion. The rest is vertebrae, cerebrospinal fluid, paraspinal tissue and air.
Yet the model runs on the \emph{whole} volume, spending almost all of its compute
on background, not on the cord. Speed matters because the segmentation is not an
end in itself: it yields quantitative biomarkers --- cross-sectional area (CSA)
and its change over time, lesion and intramedullary-tumour load --- used to track
MS, degenerative cervical myelopathy and SC injury \cite{bedard2025media,
karthik2026imag}. A segmenter too slow for the point of care limits how readily
these biomarkers reach the clinic.

\paragraph{Two routes to acceleration.}
Which strategies can make such a model faster follows directly from what
determines its cost. To first order, inference time is proportional to the number
of arithmetic operations the network performs, and for a fully-convolutional
segmentation network that count factorises into two independent terms. A 3D
convolution layer costs $\mathrm{FLOPs}\propto D\,H\,W\cdot k^3\,C_{\mathrm{in}}
C_{\mathrm{out}}$, i.e.\ \emph{linear in the number of voxels} $D H W$; summed over
the network the total is therefore
\begin{equation}\label{eq:flops}
  T \;\propto\; \mathrm{FLOPs} \;\approx\; \underbrace{N}_{\text{input size}}\;\times\;
  \underbrace{c_{\mathrm{arch}}}_{\text{architecture}},
\end{equation}
where $N$ is the number of voxels processed and $c_{\mathrm{arch}}$ is a per-voxel
cost fixed by the architecture (its depth, kernel sizes and channel widths). This
holds in the compute-bound regime typical of 3D CNNs on CPU. Because the cost is a
\emph{product} of these two factors, there are exactly two ways to reduce it.
(i)~Reduce $c_{\mathrm{arch}}$ --- replace the network with a lighter one. But
$c_{\mathrm{arch}}$ is also what buys accuracy: lighter SC and medical
segmentation architectures have so far traded accuracy for speed \cite{}, and the
field has converged on nnUNet \cite{isensee2021nnunet} precisely because it is the
most accurate. (ii)~Reduce $N$ --- crop the image to a region of interest (ROI)
around the target before segmentation, leaving the architecture, and hence its
accuracy, untouched. This \emph{localise-then-segment} strategy is already used
across many segmentation pipelines \cite{lemay2021tumor, gros2021tumor,
isensee2021nnunet, bencevic2023sts}: cropping shrinks $N$ by a factor
$\eta = N_{\mathrm{full}}/N_{\mathrm{crop}}$ (\S\ref{sec:method}) and cuts time by
the same factor per Eq.~\eqref{eq:flops}, while also removing background that would
otherwise invite false positives. To keep nnUNet's accuracy, we take route~(ii).

\paragraph{The gap: no reusable, model-agnostic crop step.}
The difficulty is that in the SC literature this localisation stage is almost
always rebuilt inside each segmentation pipeline. A dedicated network, often a
heavy 3D U-Net, is trained to localise the cord, and the crop then lives only
within that one model's loop \cite{gros2019deepseg, lemay2021tumor, gros2021tumor}.
Two consequences follow. First, every team re-implements its own cord localiser,
which duplicates effort and prevents a fair, shared baseline. Second, the speed of
localise-then-segment comes from a \emph{fast single-stage} detector (YOLO-type),
but such detectors struggle with small objects \cite{sobek2024medyolo}. So running
one directly on a small target --- a tumour, an MS lesion, oedema --- to obtain the
crop is unsafe: the target is small and variable, the detector can miss it, and
anything missed at detection is cropped away and never reaches the segmenter. The
cord is different. It is large and salient, so a fast detector localises it
reliably, and a margin-padded box around it still contains the intramedullary
pathology. That is exactly what speeds up and stabilises downstream segmentation
\cite{lemay2021tumor}. But to get that crop today, each group must first solve cord
localisation itself. What is missing is a fast, contrast-agnostic cord-cropping
step --- decoupled from any segmentation model, droppable unchanged in front of any
SC pipeline --- so that teams building new segmenters focus on training data and
segmentation, not on re-solving cord localisation.

\paragraph{Our approach.}
We fill this gap with \textbf{sc-crop}, a standalone, learned SC cropper that turns
localise-then-segment into a reusable preprocessing primitive runnable on CPU. A 2D
axial YOLO detector \cite{} runs per slice; a connected-component classifier then
removes superior false positives (brain stem, neck), yielding a tight 3D bounding
box in milliseconds (Fig.~\ref{fig:pipeline}). sc-crop replaces the heavy 3D
localiser U-Net of prior cascades with a fast 2D detector, needs no prior
whole-volume segmentation, and exposes atomic
\texttt{detect}/\texttt{crop}/\texttt{uncrop} operations so it embeds into nnUNet
preprocessing, SCT, or any custom loop. In effect, it pairs the millisecond speed
of single-stage detectors with the accuracy of a state-of-the-art nnUNet
segmenter, left unchanged. And because it shrinks the volume the segmenter sees, it
makes SC segmentation fast enough to run on CPU at the point of care. We validate
it through a controlled ablation: the same contrast-agnostic nnUNet is retrained
on cropped versus full volumes --- identical data, labels and train/val/test split,
so the crop is the only difference --- and evaluated on a frozen, multi-site,
multi-contrast test set ($n=510$), isolating the effect of cropping on accuracy.

\paragraph{Contributions.}
\begin{itemize}
  \item \textbf{C1.} A reusable, contrast- and modality-agnostic cord
        \emph{detect-and-crop} tool (2D axial YOLO + connected-component
        classifier). It ships as a CLI and Python package, with ONNX (CPU, no
        deep-learning runtime dependency) and PyTorch (GPU) backends
        \WRITE{[cervical + lumbar; MRI and CT — state scope, \S\ref{sec:exp}]}.
        This decoupled crop primitive is the methodological core.
  \item \textbf{C2.} Evidence that detect-then-segment cuts inference cost
        \WRITE{[$\times$X, EXP-SPEED]} while \emph{improving} accuracy --- with the
        nnUNet architecture left untouched. In a controlled ablation (the same
        nnUNet retrained on cropped vs full volumes; identical data, labels and
        split, so only the crop differs), cropping gives a small but highly
        significant Dice gain ($0.958$ vs $0.955$, $n=510$, paired Wilcoxon
        $p=1\times10^{-14}$), lowers variance, removes catastrophic failures
        (worst-case Dice $0.77$ vs $0.51$), and loses no ground-truth voxel
        ($510/510$).
  \item \textbf{C3.} An open-source release integrated into SCT. It includes an
        automated workflow that monitors morphometric (CSA) drift across model
        releases for lifelong learning.
\end{itemize}

% \DECISION{Page-1 figure: the pipeline overview (Fig.~\ref{fig:pipeline}) is the
% strongest opener. The empty \cite{} after "YOLO detector" needs the Ultralytics
% YOLO26 / ivadomed model_cropping reference once refs.bib is built.}

% =====================================================================
\section{Related Work}\label{sec:related}
% =====================================================================
% NOTE: literature distilled from two AI-assisted reviews (Gemini + ChatGPT).
% The Gemini review correctly frames sc-crop as a single-stage learned detector;
% the ChatGPT review wrongly calls it "heuristic/unsupervised" -- IGNORE that
% framing. Real, citable references are collected in the bibliography block at
% the end of this file. \DECISION{Verify EVERY AI-sourced citation against the
% actual paper before submission -- some entries (anonymous "2023/2024" YOLO/
% Faster-R-CNN spine papers) are unverified and must be dropped or replaced.}

\paragraph{Why a region of interest at all.}
\begin{itemize}
  \item \WRITE{Extreme class imbalance: the cord is a tiny fraction of a spinal
        volume (order $\sim$0.3\% of voxels; intramedullary lesions far less),
        so training/inferring on whole volumes wastes compute and biases models
        toward background. Cropping to a tight ROI mitigates imbalance, lowers
        GPU memory, avoids out-of-memory on large anisotropic FOVs, and
        preserves the fine in-plane resolution needed for thin structures.}
\end{itemize}

\paragraph{Classical centreline / geometric cropping.}
\begin{itemize}
  \item \WRITE{PropSeg \cite{deleener2014propseg}: elliptical Hough transform +
        deformable tubular mesh propagated along the cord; fast but sensitive to
        cord--CSF contrast and fails under severe pathology (contour leakage).}
  \item \WRITE{OptiC \cite{gros2017optic}: global centreline optimisation
        (HOG+SVM axial heatmap + rostro-caudal regularisation); validated on 501
        subjects / 20 sites; robust to contrast/resolution; powers
        \texttt{sct\_get\_centerline}.}
  \item \WRITE{SCT does ship manual cropping utilities
        (\texttt{sct\_create\_mask} $\to$ \texttt{sct\_crop\_image}), but they
        are an \emph{optional, user-driven} preprocessing requiring a prior cord
        localisation --- they are NOT invoked by the automatic segmentation
        pipeline. \DECISION{Be precise: this is the manual route a user takes if
        they want a crop; it is not what the current segmentation does.}}
\end{itemize}

\paragraph{Current best practice in SCT = no cropping (the baseline).}
\begin{itemize}
  \item \WRITE{The contrast-agnostic soft-segmentation model
        \cite{bedard2025media, karthik2026imag}, deployed as
        \texttt{sct\_deepseg spinalcord}, is the current state of the art for SC
        segmentation in SCT. It is a single nnUNet \cite{isensee2021nnunet}
        applied \emph{directly to the whole volume} --- it dropped the earlier
        localise-then-segment cascade and performs NO cropping. It is precisely
        the kind of pipeline that lacks a fast crop step in front, and motivates
        the no-crop arm of our controlled ablation.}
\end{itemize}

\paragraph{Cascade / coarse-to-fine deep learning.}
\begin{itemize}
  \item \WRITE{Sequential dual-CNN \cite{gros2019deepseg} (the older
        \texttt{sct\_deepseg\_sc}): 2D dilated-conv centreline heatmap $\to$ 3D
        patches $\to$ 3D CNN segmentation; median Dice $\sim$95\% for the cord.
        Its localisation step is effectively a learned crop --- but this cascade
        has since been \emph{superseded} in SCT by the direct, no-crop
        contrast-agnostic model above. \DECISION{Distinguish clearly:
        \texttt{sct\_deepseg\_sc} (old, cascade) vs \texttt{sct\_deepseg
        spinalcord} (current, contrast-agnostic, full-volume, no crop) --- they
        are different methods; do not conflate.}}
  \item \WRITE{Locate-and-label \cite{gros2021tumor}: a light 3D U-Net emits a
        tight 3D bbox, multi-contrast images are cropped to it, then a
        multi-class 3D U-Net segments intramedullary tumours/oedema/cavities.
        Closest in spirit to detect-then-segment, but uses a heavy 3D U-Net to
        localise.}
  \item \WRITE{Lemay et al.\ \cite{lemay2021tumor} --- KEY PRECEDENT: 3D U-Net
        localises the cord, the volume is cropped to the mask bbox, then tumour
        segmentation runs on the crop. Reports cord localisation Dice 88.7\%
        (100\% detection), tumour Dice +5\,pt and oedema +30\,pt vs.\ training on
        the full image, and a faster pipeline (22\,s vs.\ 48\,s without the crop
        step). This is the canonical evidence that cropping speeds up inference
        AND improves Dice --- position sc-crop directly against it.}
  \item \WRITE{nnU-Net low/high-res 3D cascade \cite{isensee2021nnunet}:
        auto-triggered when volumes greatly exceed the 3D patch size; a coarse
        3D U-Net provides global context to a fine one. Powerful but the coarse
        stage still runs a full 3D network over the (downsampled) whole volume.}
  \item \WRITE{Coarse-to-fine ROI cascades in other organs (pancreas, small
        liver tumours): \cite{zhu2018c2f, zhou2018pancreas, wu2023liver}.}
  \item \WRITE{Segment-then-Segment \cite{bencevic2023sts}: coarse low-res
        segmentation $\to$ salient high-res crops $\to$ fine per-crop
        segmentation; improves small-structure recall vs.\ plain downsampling.}
  \item \WRITE{Limitation of axis-aligned bounding boxes (AABB)
        \cite{shin2024aabb}: AABBs over-cover elongated/curved structures and
        propagate first-stage false negatives --- motivates a tight, cord-aware
        box and a safety margin (our per-face mm padding + \texttt{check\_label\_crop}).}
\end{itemize}

\paragraph{Single-stage detectors (the sc-crop family).}
\begin{itemize}
  \item \WRITE{YOLO-family single-pass detectors give millisecond-scale boxes;
        widely used in spine imaging (vertebral-level/stenosis detection, disc
        pathology). Run on 2D slices, they suit the cord's tubular geometry and
        anisotropic acquisition while staying cheap on CPU. ivadomed's
        \texttt{model\_cropping\_sc\_contrast-agnostic\_yolo} trains such a
        detector for contrast-agnostic cord cropping --- the detector sc-crop
        packages.}
  \item \WRITE{Detection frameworks for context: nnDetection (YOLO-style heads
        on nnU-Net), MedYOLO (3D), MONAI/TorchIO crop transforms --- mostly
        tied to a training loop or toolbox, not a standalone reusable crop.}
\end{itemize}

\paragraph{Positioning of sc-crop (closing paragraph).}
\begin{itemize}
  \item \WRITE{sc-crop is a \emph{learned} single-stage detector (2D axial
        YOLO26n + a connected-component classifier regulariser), NOT a heuristic
        centred crop. Unlike PropSeg/OptiC chains it needs no prior whole-volume
        segmentation and no multi-tool sequence; unlike Lemay/Locate-and-label it
        replaces the heavy 3D localiser U-Net with a fast 2D detector; unlike
        nnDetection/MONAI it is \emph{decoupled} from any training loop or
        toolbox. It is distributed as ONNX (CPU, no torch/numpy conflicts) with
        a PyTorch/GPU path, is contrast- and modality-agnostic (MRI + CT,
        cervical + lumbar), and exposes atomic \texttt{detect/crop/uncrop} so it
        drops into nnU-Net preprocessing, SCT, or any custom pipeline.}
  \item \DECISION{Reuse the comparison table from sc-crop ARCHITECTURE.md
        (MONAI/TorchIO/SCT/nnDetection/HF/Ultralytics: interface, model hosting,
        offline-vs-online crop) as Related-Work Table or an appendix.}
\end{itemize}

% =====================================================================
\section{Method}\label{sec:method}
% =====================================================================
\begin{itemize}
  \item \textbf{Overview (Fig.~\ref{fig:pipeline}).}
    \begin{itemize}
      \item \WRITE{detect $\rightarrow$ crop $\rightarrow$ segment $\rightarrow$
            uncrop, with the invariant: \emph{identical padding at train and
            inference time}.}
    \end{itemize}
  \item \textbf{Detector (Fig.~\ref{fig:arch}).}
    \begin{itemize}
      \item \WRITE{YOLO26n, single class ``sc'', run per axial slice. Input
            preprocessing: reorient to LAS; resample to SI=10\,mm,
            in-plane=1\,mm; 3-channel slices (superior/current/inferior
            neighbours); normalisation (\texttt{slice\_all}: 0.5/99.5
            percentile per slice). Letterbox to 320 px (rectangular,
            multiple-of-32), bit-exact with YOLO's predict().}
      \item \WRITE{Per-slice output -> keep max-confidence box per slice
            (conf $\geq$ 0.1).}
    \end{itemize}
  \item \textbf{Regularization (the key robustness component).}
    \begin{itemize}
      \item \WRITE{Detections are grouped into SI connected components. A second
            YOLO \emph{classifier} (``sc'' vs ``no\_sc'', square 320 letterbox)
            scans components from most-superior downward and keeps the first
            validated component and everything inferior to it — removing
            superior false positives (brain stem, neck structures).}
      \item \WRITE{Alternative \texttt{graphtrim} regularizer (geometry-only:
            SI gap $\geq$40\,mm or per-hop face shift) as a no-classifier
            fallback; \texttt{none} for ablation.}
    \end{itemize}
  \item \textbf{3D bounding box + padding.}
    \begin{itemize}
      \item \WRITE{Aggregate per-slice boxes to a 3D box in native LAS voxels;
            pad per face in mm (defaults sup=40, inf=100, L=R=15, ant=15,
            post=22; inferior is large to cover lumbar/anisotropic cords);
            clamp to image bounds; reorient box back to native orientation.}
      \item \WRITE{\texttt{check\_label\_crop} / \texttt{CropReport}: QC that no
            GT voxel is cut, reporting extra mm needed per face (how the
            defaults were tuned).}
    \end{itemize}
  \item \textbf{crop / uncrop.}
    \begin{itemize}
      \item \WRITE{crop updates the affine so the crop sits at the correct
            world position; uncrop pastes the segmentation back into the full
            original space/affine.}
    \end{itemize}
  \item \textbf{Computational impact of cropping (backend-independent argument).}
    \begin{itemize}
      \item \WRITE{This instantiates route~(ii) of Eq.~\eqref{eq:flops}: with
            $c_{\mathrm{arch}}$ fixed (the nnUNet is unchanged), inference cost is
            linear in the voxel count $N$. A 3D conv layer costs FLOPs
            $\propto D\,H\,W\,k^3 C_{in}C_{out}$, so cropping the input from
            $D\times H\times W$ to $d\times h\times w$ (with margin $m$) reduces
            compute by $\eta = N_{\mathrm{full}}/N_{\mathrm{crop}} =
            \frac{D H W}{d h w}$, typically $\eta\!\approx\!5$--$15$ for localised
            spinal cords. This is the speed argument \emph{attributable to cropping
            alone}, independent of inference backend (PyTorch/ONNX) --- use it to
            motivate the measured speedup and to decouple the crop effect from the
            ONNX-export effect.}
      \item \DECISION{Report the empirical FOV-reduction factor $\eta$ measured
            on the test set (crop voxels / original voxels) alongside the
            theoretical bound --- this is the clean, confound-free efficiency
            result.}
    \end{itemize}
  \item \textbf{Inference backends \& engineering for reuse.}
    \begin{itemize}
      \item \WRITE{ONNX Runtime (CPU, no ultralytics/torch at runtime) vs
            PyTorch (.pt, GPU). SHA256-pinned, versioned model cache;
            \texttt{SC\_CROP\_CACHE\_DIR} for SCT/air-gapped HPC.}
      \item \WRITE{Atomic public API: detect() / crop() / uncrop() — what makes
            it embeddable (nnUNet preprocessing, SCT, custom loops).}
    \end{itemize}
  \item \textbf{Design choices to defend explicitly (Discussion ties back).}
    \begin{itemize}
      \item \WRITE{Offline cropping (never in the training loop) — matches
            nnDetection/MONAI/SCT practice; YOLO too slow per-epoch.}
      \item \WRITE{2D-slice detection vs 3D detector — speed, robustness to
            anisotropy, data efficiency.}
      \item \WRITE{ONNX distribution — eliminates dependency conflicts.}
    \end{itemize}
  \item \textbf{Segmentation model trained on crops (C2).}
    \begin{itemize}
      \item \WRITE{nnUNet 3d\_fullres; training data = sc-crop-cropped
            image/label pairs (same padding as inference); GT cleaned to
            largest connected component. Reference the reproducible script
            (\texttt{train\_contrast\_agnostic.sh}, 5 steps).}
      \item \FUTURE{If a custom nnUNet trainer that simulates random YOLO-style
            crops during training is added, describe it here (augmentation:
            GT bbox + 0--10\,mm / 10\,mm--max padding, 50/50). Reserve a
            Method subsection for it.}
    \end{itemize}
\end{itemize}

% =====================================================================
\section{Experiments}\label{sec:exp}
% =====================================================================
\begin{itemize}
  \item \textbf{Datasets \& contrasts (Table~\ref{tab:datasets}).}
    \begin{itemize}
      \item \RESULT{14 datasets, 9 contrasts, frozen held-out test set
            $n=510$. Per-contrast $n$ (from \texttt{metrics.json}):
            T2w=159, T2star=94, T1w=56, MTS flip-1/flip-2 = 49/49,
            DWI(rec-average)=49, PSIR=33, UNIT1(MP2RAGE)=11, STIR=10.}
      \item \RESULT{Datasets present: data-multi-subject (294), canproco (85),
            sct-testing-large (36), PSIR/dcm-brno (28), dcm-zurich (14),
            sci-zurich (14), basel-mp2rage (11), sci-colorado (8),
            lumbar-vanderbilt (6), dcm-zurich-lesions-20231115 (5),
            lumbar-epfl (3), sci-paris (2), dcm-zurich-lesions (2),
            site\_006/007 (1/1).}
      \item \WRITE{Answer the modality question explicitly: MRI multi-contrast
            (T1, T2, T2*, MT, DWI, PSIR, STIR, MP2RAGE); sc-crop \emph{detector}
            also supports CT (auto HU background detection) — state whether CT
            is evaluated here or only claimed/qualitative.}
      \item \WRITE{Frozen split rationale + reproducibility (predefined
            datasplits with pinned dataset version commits).}
    \end{itemize}
  \item \textbf{Metrics.}
    \begin{itemize}
      \item \WRITE{Segmentation: Dice (soft/binary — state which) vs cleaned GT.}
      \item \WRITE{Crop quality: fraction of subjects with zero GT voxels lost
            (\texttt{check\_label\_crop}); FOV reduction factor.}
      \item \WRITE{Efficiency: wall-clock inference time per volume, CPU and
            GPU; peak memory; report hardware.}
      \item \WRITE{Morphometric agreement: C2--C3 CSA vs reference; CSA std
            across contrasts (lifelong-drift metric).}
      \item \FUTURE{Add detector-only metrics if you run a detection ablation
            (recall of cord coverage, bbox IoU) — optional.}
    \end{itemize}
  \item \textbf{Baselines.}
    \begin{itemize}
      \item \WRITE{Primary comparison = controlled crop ablation: the same
            contrast-agnostic nnUNet trained on FULL volumes (no crop) vs on
            sc-crop crops, with identical data, labels and train/val/test split.
            Only the crop differs, so any Dice difference is attributable to
            cropping. (The previously deployed model is intentionally NOT used as
            an accuracy baseline: it was under-trained, so a comparison against it
            would not isolate cropping.)}
      \item \WRITE{Pipeline variants for the speed comparison: sc-crop + nnUNet
            PyTorch (\texttt{sc-segment-pt}) and ONNX (\texttt{sc-segment-onnx}).}
      \item \FUTURE{Reserve baseline slots for future model versions (vX.Y).}
    \end{itemize}
  \item \textbf{Implementation / reproducibility.}
    \begin{itemize}
      \item \WRITE{Versions: sc-crop v0.4.1, detector model v0.0.10, nnUNet
            v2.6.0 pin, PyTorch 2.8; cite VERSIONS.md mapping. Release the
            frozen split YAML and the training/test shell scripts.}
    \end{itemize}
\end{itemize}

% =====================================================================
\section{Results}\label{sec:results}
% =====================================================================
\begin{itemize}
  \item \textbf{R1 — Segmentation accuracy: crop ablation (Table~\ref{tab:dice}, done).}
    \begin{itemize}
      \item \RESULT{Controlled ablation, crop vs no-crop ($n=510$): cropping
            raises overall Dice from $0.9548\pm0.0282$ (no-crop) to
            $0.9578\pm0.0201$ (crop). Paired Wilcoxon $p=1.1\times10^{-14}$;
            cropping wins in $64.9\%$ of subjects (median $\Delta=+0.0019$).
            Source: \texttt{qc\_results\_dataset7000} (crop) and
            \texttt{qc\_results\_dataset7001} (no-crop).}
      \item \RESULT{Robustness is the larger effect: cropping cuts variance by
            $\sim$29\% ($\sigma=0.0201$ vs $0.0282$), raises the worst case from
            Dice $0.51$ to $0.77$, and reduces sub-$0.90$ subjects ($9$ vs $11$).}
      \item \RESULT{Per-contrast (Table~\ref{tab:dice}): cropping improves mean
            Dice in \emph{all 9} contrasts and never degrades any. The gain is
            significant for T2w, T2$^\ast$, T1w, MTS (both flips) and STIR, and at
            parity on the well-covered/small-$n$ contrasts (DWI, PSIR, UNIT1).}
      \item \WRITE{Ready-to-use sentence (EN): ``Retraining the same
            contrast-agnostic nnUNet on cropped instead of full volumes improves
            mean Dice in all nine contrasts (significant overall, paired Wilcoxon
            $p=1.1\times10^{-14}$), and its main benefit is robustness: it lowers
            variance and removes catastrophic failures (worst-case Dice $0.51
            \rightarrow 0.77$).'' (FR): ``R\'eentra\^iner le m\^eme mod\`ele sur
            des volumes crop\'es plut\^ot que complets am\'eliore le Dice moyen
            dans les neuf contrastes (significatif globalement, $p=1.1\times
            10^{-14}$), et son b\'en\'efice principal est la robustesse : variance
            r\'eduite et disparition des \'echecs catastrophiques.''}
      \item \FUTURE{Empty column reserved for vX.Y / future models in
            Table~\ref{tab:dice}.}
    \end{itemize}
  \item \textbf{R2 — Crop quality (Table~\ref{tab:crop}).}
    \begin{itemize}
      \item \RESULT{$0/510$ subjects lost any GT voxel (n\_bad\_crops $=0$).
            Source: \texttt{metrics.json} aggregate.}
      \item \WRITE{Report FOV reduction (mean voxels\_after / voxels\_before, or
            volume ratio) to quantify how much background is removed.}
    \end{itemize}
  \item \textbf{R3 — Inference speed (Table~\ref{tab:speed}).}
    \begin{itemize}
      \item \DECISION{Drop ONNX from the speed comparison. The issue-\#2
            ``$\sim$10$\times$'' is \texttt{sc-segment-onnx} vs
            \texttt{sct\_deepseg} --- it mixes the ONNX export with the crop, so
            it cannot show that cropping is what speeds things up. Compare
            SAME-backend instead.}
      \item \WRITE{Empirical FOV reduction $\eta$ = crop voxels / original
            voxels on the test set (the clean, backend-free efficiency number);
            report mean and range, expected $\approx$5--15.}
      \item \WRITE{Like-for-like wall-clock: \texttt{sct\_deepseg} (full volume)
            vs \texttt{sc-segment-pt} (crop + nnUNet, PyTorch), same machine;
            split detector vs segmentation vs total, CPU and GPU. \DECISION{The
            cleanest version measures the SAME nnUNet on full vs cropped input
            --- ties to the optional no-crop ablation in \S\ref{sec:exp}.}}
      \item \WRITE{Keep ONNX only as a separate ``deployment'' note (no torch at
            runtime, $\sim$10$\times$ vs baseline), clearly NOT attributed to
            cropping.}
      \item \FUTURE{Reserve rows for future backends/models.}
    \end{itemize}
  \item \textbf{R4 — Morphometric / lifelong monitoring (Fig.~\ref{fig:csa}).}
    \begin{itemize}
      \item \WRITE{CSA C2--C3 violin plots per contrast across model releases;
            CSA std across contrasts as the drift metric. These are produced
            automatically by the GitHub Action on each release
            (\texttt{morphometric\_plots.zip}).}
      \item \FUTURE{This figure literally grows by one model per release — it
            is the natural ``future models'' result. Describe it as a living
            benchmark.}
    \end{itemize}
  \item \textbf{R5 — Qualitative (Fig.~\ref{fig:qual}).}
    \begin{itemize}
      \item \WRITE{Crops + segmentations across the hardest contrasts
            (PSIR/STIR/MP2RAGE, lumbar, SCI lesions); show a failure/edge case
            honestly.}
    \end{itemize}
\end{itemize}

% =====================================================================
\section{Discussion and Limitations}\label{sec:discussion}
% =====================================================================
\begin{itemize}
  \item \textbf{Discussion.}
    \begin{itemize}
      \item \WRITE{Why crop-then-segment helps: removes background false
            positives, lets nnUNet spend capacity/resolution on the cord;
            decouples a reusable crop step from the segmentation model.}
      \item \WRITE{Practical impact: $\sim$10$\times$ CPU speedup makes
            laptop/clinical-PC inference feasible; ONNX removes dependency
            friction; SCT integration = real deployment.}
      \item \WRITE{Generality: same crop tool for other SC tasks (lesions,
            tumors, gray matter) and CT.}
    \end{itemize}
  \item \textbf{Limitations.}
    \begin{itemize}
      \item \WRITE{The mean Dice gain from cropping is small ($+0.0029$, $\sim
            0.3$\,pt) though highly significant ($p=1.1\times10^{-14}$); the
            primary accuracy benefit is robustness (lower variance, removal of
            catastrophic failures) rather than a large mean jump. Because the
            ablation is fully controlled --- only the crop differs --- the gain is
            attributable to cropping, but its magnitude should not be oversold.}
      \item \WRITE{Test set is imbalanced across contrasts/pathologies
            (T2w-heavy; STIR/MP2RAGE small $n$); few severe-pathology cases.}
      \item \WRITE{Fixed padding assumes the detector+margin always contains
            the cord — quantify residual risk (here $0/510$, but acknowledge
            distribution shift / extreme FOVs).}
      \item \WRITE{CT support is implemented but report its actual evaluation
            scope honestly.}
      \item \WRITE{Detector failure modes: very anisotropic sagittal,
            superior false positives handled by cls but not infallible.}
    \end{itemize}
\end{itemize}

% =====================================================================
\section{Conclusion}\label{sec:conclusion}
% =====================================================================
\begin{itemize}
  \item \WRITE{Restate: a lightweight, contrast-agnostic, reusable detect+crop
        step improves SC segmentation accuracy and substantially cuts
        inference cost, is deployed in SCT, and is monitored over time.}
  \item \WRITE{Outlook: more contrasts/pathologies per release; the living CSA
        benchmark; crop step as a general SC-ROI primitive.}
\end{itemize}

% =====================================================================
\section*{Data and Code Availability}
% =====================================================================
% Required by MedIA / Imaging Neuroscience.
\begin{itemize}
  \item \WRITE{sc-crop (CLI + package, ONNX/PT models, SHA256-pinned releases);
        contrast-agnostic training/eval scripts + frozen split; SCT integration
        branch; GitHub Action for CSA monitoring. List repos + versions.}
\end{itemize}

% =====================================================================
%  FIGURES & TABLES — skeletons (fill captions/data)
% =====================================================================

\begin{figure}[t]\centering
  \fbox{\parbox{0.9\linewidth}{\centering \WRITE{Fig 1 — pipeline:
  detect $\to$ crop $\to$ segment $\to$ uncrop, with ``same padding train/infer''
  invariant. Reuse the sc-crop README banner as a base.}}}
  \caption{End-to-end detect-then-segment pipeline.}
  \label{fig:pipeline}
\end{figure}

\begin{figure}[t]\centering
  \fbox{\parbox{0.9\linewidth}{\centering \WRITE{Fig 2 — sc-crop architecture:
  2D axial YOLO detector per slice + connected-component classifier
  regularization + 3D bbox aggregation + mm padding. The debug-panel mosaic
  (\texttt{save\_debug\_panel}) makes a great real example.}}}
  \caption{sc-crop detector and connected-component regularization.}
  \label{fig:arch}
\end{figure}

\begin{figure}[t]\centering
  \fbox{\parbox{0.9\linewidth}{\centering \FUTURE{Fig 3 (CSA monitoring) —
  reuse \texttt{imag.a.1105\_fig2.png} / the auto-generated
  \texttt{morphometric\_plots.zip}; grows by one model per release.}}}
  \caption{Lifelong CSA-drift monitoring across model releases.}
  \label{fig:csa}
\end{figure}

\begin{figure}[t]\centering
  \fbox{\parbox{0.9\linewidth}{\centering \WRITE{Fig 4 — qualitative crops +
  segmentations across hard contrasts; include one honest failure case.}}}
  \caption{Qualitative results across contrasts.}
  \label{fig:qual}
\end{figure}

\begin{table}[t]\centering
  \caption{Test set: datasets and contrasts ($n=510$). \RESULT{numbers from
  \texttt{metrics.json}}.}
  \label{tab:datasets}
  \begin{tabular}{lrr}
    \toprule
    Contrast & $n$ & \WRITE{Datasets} \\
    \midrule
    T2w        & 159 & \WRITE{} \\
    T2star     & 94  & \WRITE{} \\
    T1w        & 56  & \WRITE{} \\
    MTS flip-1 & 49  & \WRITE{} \\
    MTS flip-2 & 49  & \WRITE{} \\
    DWI        & 49  & \WRITE{} \\
    PSIR       & 33  & \WRITE{} \\
    MP2RAGE (UNIT1) & 11 & \WRITE{} \\
    STIR       & 10  & \WRITE{} \\
    \midrule
    \textbf{Total} & \textbf{510} & 14 datasets \\
    \bottomrule
  \end{tabular}
\end{table}

\begin{table}[t]\centering
  \caption{Crop ablation: segmentation accuracy by contrast (Dice, mean$\pm$std).
  no-crop = the contrast-agnostic nnUNet trained on FULL volumes; crop = the same
  nnUNet trained on sc-crop crops (identical data, labels and split --- only the
  crop differs). $\Delta$ = crop $-$ no-crop; $p$ = paired Wilcoxon per contrast.
  \RESULT{All values computed from \texttt{qc\_results\_dataset7001} (no-crop) and
  \texttt{qc\_results\_dataset7000} (crop), $n=510$.}}
  \label{tab:dice}
  \begin{tabular}{lrcccc}
    \toprule
    Contrast & $n$ & no-crop & crop & $\Delta$ & $p$ \\
    \midrule
    T2w               & 159 & $0.9537\pm0.0428$ & $0.9570\pm0.0250$ & $+0.0033$ & $1.9\times10^{-2}$ \\
    T2star            & 94  & $0.9615\pm0.0136$ & $0.9649\pm0.0111$ & $+0.0033$ & $2.3\times10^{-5}$ \\
    T1w               & 56  & $0.9548\pm0.0193$ & $0.9581\pm0.0181$ & $+0.0033$ & $7.4\times10^{-6}$ \\
    MTS flip-1        & 49  & $0.9586\pm0.0156$ & $0.9622\pm0.0121$ & $+0.0036$ & $8.8\times10^{-4}$ \\
    MTS flip-2        & 49  & $0.9510\pm0.0249$ & $0.9545\pm0.0259$ & $+0.0035$ & $7.7\times10^{-4}$ \\
    DWI               & 49  & $0.9570\pm0.0092$ & $0.9574\pm0.0094$ & $+0.0005$ & $0.49$ (ns) \\
    PSIR              & 33  & $0.9406\pm0.0190$ & $0.9413\pm0.0195$ & $+0.0007$ & $0.83$ (ns) \\
    UNIT1 (MP2RAGE)   & 11  & $0.9616\pm0.0093$ & $0.9646\pm0.0113$ & $+0.0030$ & $0.17$ (ns) \\
    STIR              & 10  & $0.9380\pm0.0103$ & $\mathbf{0.9434\pm0.0089}$ & $+0.0054$ & $3.9\times10^{-3}$ \\
    \midrule
    \textbf{Overall}  & 510 & $0.9548\pm0.0282$ & $\mathbf{0.9578\pm0.0201}$ & $+0.0029$ & $1.1\times10^{-14}$ \\
    \bottomrule
  \end{tabular}
  \\[2pt]\footnotesize Overall paired Wilcoxon (crop vs no-crop, $n=510$):
  $p=1.1\times10^{-14}$, median $\Delta=+0.0019$, crop wins $64.9\%$. Cropping
  improves mean Dice in all 9 contrasts and never degrades any; the variance drop
  ($\sigma=0.0201$ vs $0.0282$) and worst-case improvement ($0.51\rightarrow0.77$)
  are the dominant effect.
\end{table}

\begin{table}[t]\centering
  \caption{Crop quality (\texttt{check\_label\_crop}).}
  \label{tab:crop}
  \begin{tabular}{lc}
    \toprule
    Metric & Value \\
    \midrule
    Subjects with zero GT voxel lost            & $510/510$ \RESULT{(n\_bad\_crops=0)} \\
    Mean label-voxel preservation ratio         & $1.0000$ (min $1.0000$) \RESULT{} \\
    Mean FOV reduction (crop vol / orig vol)     & \WRITE{NOT in metrics.json --- compute separately} \\
    \bottomrule
  \end{tabular}
  \\[2pt]\footnotesize \DECISION{In \texttt{metrics.json},
  \texttt{voxels\_before}/\texttt{voxels\_after} are \emph{label} voxels
  (crop-safety: total vs kept), NOT field of view. They are equal for every
  subject ($\Rightarrow$ 0 lost). To report how much background the crop
  removes (a key efficiency argument), compute cropped-volume/original-volume
  separately from the bbox vs image dimensions.}
\end{table}

\begin{table}[t]\centering
  \caption{Inference speed per volume, SAME backend on both sides so the
  difference is attributable to cropping (not to ONNX). \WRITE{Fill from a
  controlled timing run; add a GPU block.}}
  \label{tab:speed}
  \begin{tabular}{lcccc}
    \toprule
    Method (PyTorch backend) & FOV factor $\eta$ & Detect (s) & Segment (s) & Total (s) \\
    \midrule
    \texttt{sct\_deepseg spinalcord} (full volume) & $1$ & --- & \WRITE{} & \WRITE{} \\
    \texttt{sc-segment-pt} (crop + nnUNet)         & \WRITE{$\approx$5--15} & \WRITE{} & \WRITE{} & \WRITE{} \\
    \FUTURE{future model}                          &     &     &     &     \\
    \bottomrule
  \end{tabular}
  \\[2pt]\footnotesize \DECISION{ONNX (\texttt{sc-segment-onnx}) is deliberately
  excluded here --- its $\sim$10$\times$ vs baseline mixes the ONNX export with
  the crop. Mention ONNX only as a deployment note in the text.}
\end{table}

% =====================================================================
\section*{Authoring checklist (delete before submission)}
% =====================================================================
\begin{itemize}
  \item \DECISION{Run the no-crop nnUNet ablation — biggest reviewer risk.}
  \item \DECISION{Get exact speed numbers SAME-backend (crop vs full, CPU+GPU)
        into Tab.~\ref{tab:speed}; ONNX excluded from the attribution.}
  \item \WRITE{Compute empirical FOV factor $\eta$ on the test set.}
  \item \DECISION{Verify EVERY AI-sourced citation below against the real paper;
        drop the unverified anonymous spine-YOLO/Faster-R-CNN entries.}
  \item \WRITE{Build \texttt{refs.bib} from the keys used in Related Work and add
        \texttt{\textbackslash bibliography\{refs\}}.}
  \item \DECISION{State CT evaluation scope (evaluated vs claimed).}
  \item \DECISION{Confirm author list + target venue + template.}
\end{itemize}

% =====================================================================
%  REFERENCES — seed list for refs.bib (VERIFY each before use)
% =====================================================================
% Verified, clearly-real (use these):
%   deleener2014propseg  De Leener et al., "Robust, accurate and fast automatic
%                        segmentation of the spinal cord" (PropSeg), NeuroImage 2014.
%   gros2017optic        Gros et al., "Automatic spinal cord localization,
%                        robust to MRI contrasts using global curve optimization"
%                        (OptiC), Medical Image Analysis 2018 (501 subj / 20 sites).
%   gros2019deepseg      Gros et al., "Automatic segmentation of the spinal cord
%                        and intramedullary MS lesions on MRI" (sct_deepseg_sc),
%                        NeuroImage 2019.
%   gros2021tumor        Gros et al., locate-and-label intramedullary tumour
%                        segmentation (multi-class 3D U-Net cascade).
%   lemay2021tumor       Lemay et al., "Automatic Multiclass Intramedullary
%                        Spinal Cord Tumor Segmentation on MRI with Deep
%                        Learning", NeuroImage:Clinical 31:102766, 2021.
%                        (KEY precedent: crop 22s vs 48s, Dice +5pt tumour.)
%   isensee2021nnunet    Isensee et al., "nnU-Net", Nature Methods 2021.
%   bencevic2023sts      Benčević et al., "Segment-then-Segment", Sensors
%                        23(2):633, 2023.
%   cheng2021vertebrae   Cheng et al., two-stage Dense-U-Net vertebrae
%                        localization+segmentation, Sci. Rep. 11:22156, 2021.
%   zhu2018c2f           Zhu et al., "A 3D Coarse-to-Fine Framework for
%                        Volumetric Medical Image Segmentation", 3DV 2018.
%   zhou2018pancreas     Zhou et al., dense-regression pancreas, MICCAI 2018.
%   wu2023liver          Wu et al., coarse-to-fine small liver tumour
%                        detection+segmentation, Diagnostics 13(15):2504, 2023.
%   shin2024aabb         Shin et al. (2024), limits of axis-aligned bounding
%                        boxes for elongated structures. [VERIFY exact ref]
%   bedard2025media      Bédard, Karthik et al., "Towards contrast-agnostic soft
%                        segmentation of the spinal cord", Medical Image
%                        Analysis 101:103473, 2025. (already in repo README)
%   karthik2026imag      Karthik et al., "Monitoring morphometric drift in
%                        lifelong learning segmentation of the spinal cord",
%                        Imaging Neuroscience 2026. (repo README)
%   cohenadad2021spinegeneric  Cohen-Adad et al., Spine-Generic open-access
%                        qMRI dataset, Sci. Data 8:219, 2021.
%   sobek2024medyolo     Sobek et al., "MedYOLO: A Medical Image Object Detection
%                        Framework", J. Imaging Inform. Med. 2024. (3D single-stage
%                        detector; cited for small-object difficulty in Intro §gap.)
%                        [VERIFY exact venue/year before submission]
%
% UNVERIFIED (do NOT cite until confirmed; came from AI reviews as "anonymous"):
%   - "YOLO11 cervical-compression 2024", "ViT axial-GRE compression 2024",
%     "Faster R-CNN SCI 2023", "YOLOv4 spinal neoplasms 2023". Either find the
%     real papers or drop them.
%
% TOOLS (cite software/manuals as needed): SCT, MONAI, TorchIO, nnDetection,
%   Ultralytics YOLO, ivadomed/model_cropping_sc_contrast-agnostic_yolo.

\end{document}
